\section{How large is practically feasible?}
\label{sec:scale}

A useful specification is $(k,n,\alpha,\text{ensemble},\text{task},
\text{time and hardware budget},\text{success target})$.
``A million variables'' alone says little. A sparse satisfiable formula
with a readily discoverable assignment and a near-critical UNSAT formula
can have radically different costs. General industrial SAT results also
cannot be read as uniform-random-$3$-SAT capacity estimates.

\subsection{Concrete evidence, including failures}
The historical Configurable SAT Solver Challenge used mixed random $3$-CNF
instances with 350 variables and 1493 clauses, and a satisfiable family with
1000 variables and 4260 clauses \citep[Appendix A.3]{hutter2017}. These
sizes indicate how small challenging complete-decision benchmarks can be;
they are not present-day hardware ceilings.

An especially informative comparison is the deliberately difficult
\texttt{Large\_Threshold} benchmark in \citet[Table 2]{luo2014}.
Table~\ref{tab:hardinstances} preserves the low success rates rather than
reporting only the largest instance size attempted.
\begin{table}[htbp]
\centering\small
\begin{tabular}{@{}rrrrr@{}}
\toprule
$k$ & $n$ & Density & probSAT solved & DCCASat solved \\
\midrule
3 & 15\,000 & 4.267 & 2/100 & 6/100 \\
4 & 2\,500 & 9.931 & 5/100 & 5/100 \\
5 & 600 & 21.117 & 2/100 & 9/100 \\
6 & 300 & 43.37 & 0/100 & 11/100 \\
7 & 150 & 87.79 & 28/100 & 43/100 \\
\bottomrule
\end{tabular}
\caption{Published 2014 local-search measurements: one run per instance,
5000-second cutoff, GNU/Linux on an Intel Core i7 at 2.4 GHz
(two cores reported in the setup). These are counts of witnesses found,
not proofs that the remaining formulas are UNSAT. The original benchmark
does not establish that every generated instance is satisfiable.}
\label{tab:hardinstances}
\end{table}

In contrast, \citet[Figures 3--5]{marino2016} report BSP solutions for
random $3$-SAT with $n=10^6$ at densities close to, but below, the predicted
transition, including analysis at $\alpha=4.262$. This is evidence that
million-variable \emph{search} can be feasible with a specialized method.
It is not a certificate of efficient complete decision at the transition.
The extrapolated value near $4.268$ in that paper is obtained by fitting
residual complexity; it must not be quoted as successful solving of
million-variable instances at that density. The cited figures also do not
provide a matched wall-clock and hardware comparison with Table~\ref{tab:hardinstances}.

The useful planning conclusion is therefore conditional: expect useful
near-critical complete-solver studies to begin at hundreds of variables;
thousands may already demand substantial budgets. For witness search below
the transition, thousands to hundreds of thousands are natural scales to
test, and a million is demonstrated in specialized published experiments.
These are starting expectations, not a machine-independent boundary or a
promise that every density in those ranges is tractable. Larger $k$ can be
hard at much smaller $n$, as the table illustrates.

\subsection{Input size and memory also constrain experiments}
There are $km=k\alpha n$ literal occurrences. Storing only one 32-bit signed
index per occurrence costs about $4k\alpha n$ bytes. At representative
densities this gives:
\begin{center}\small
\begin{tabular}{@{}rrrr@{}}
\toprule
$k$ & $n$ & $\alpha$ & Literal indices alone \\
\midrule
3 & $10^6$ & 4.267 & 51.2 MB \\
5 & $10^5$ & 21.117 & 42.2 MB \\
10 & $10^5$ & 708.936 & 2.84 GB \\
\bottomrule
\end{tabular}
\end{center}
These are arithmetic estimates in decimal units, not memory measurements.
Clause offsets, watches, occurrence lists, learned clauses, proof logs and
batched GPU states add overhead, sometimes substantial. Near the threshold,
$m$ grows exponentially with $k$; memory and input generation can become
limiting before the number of variables looks large.

\subsection{A reproducible first feasibility study}
The following is a \emph{proposed} pilot, not an experiment performed here.
For uniform random $3$-SAT, use densities $3.5,4.0,4.2,4.25,4.267,4.30,4.50$.
The values around $4.267$ are empirical landmarks, not a known exact
finite-$n$ threshold. For complete solving, begin with
$n=100,200,400,800$ and extend to $1600$ and beyond only after checking the
timeout rate. For witness search at the less critical densities, try
$n=10^3,10^4,10^5$ before committing to $10^6$.

Use at least 100 independently generated instances per cell if resources
permit, initially with a 60-second cutoff; repeat selected cells at 300 or
3600 seconds. These are budget choices, not asserted requirements for
success. Record generation seeds, solver versions, parameters, CPU cores,
RAM, GPU model and precision. Retain each raw result as SAT with a checked
witness, UNSAT with the stated checking status, or unknown. Report success
fractions and runtime distributions, including censored runs, rather than
only the median runtime among successes.

For $N$ sampled formulas, $S$ verified SAT and $U$ verified UNSAT outputs give
an empirical SAT-fraction interval
\begin{equation}
 [S/N,\;1-U/N].
 \label{eq:censor}
\end{equation}
This accounts for unknown answers in that sample; estimating the ensemble
probability additionally needs sampling uncertainty. Treating all timeouts
as UNSAT biases a threshold estimate downward. Successful-only datasets or
planted instances answer a different question.
